\documentclass[12pt]{article} 

% --- Packages for formatting ---
\usepackage[utf8]{inputenc} 
\usepackage[margin=1in]{geometry} 
\usepackage{graphicx} 
\usepackage{hyperref} 

% --- Document Information ---
\title{Robotics Project Report}
\author{F331073}
\date{\today}

\begin{document}

\maketitle 

\tableofcontents 
\newpage

\section{Abstract}
There are many different ways to tackle this project, all of which could
result in a successful robot. In this report, we will explore the various
approaches and methodologies that can be employed to navigate the circuit
as quickly as possible.

\section{Introduction}

\section{Related Works} 

\section{Datasets}

\subsection{Data Collection} 
There are many different approaches to collecting data for this project. Anything from
manual image collection to simulations, data can be gathered in many ways.
I chose to opt for a more manual approach where I made use of my previously defined
movement functions---using w/a/s/d---to control the robot around the track. To allow
the robot to then record the data, I made use of threading. This allowed for parallel
recording and motion operations. 

Initially, I was unsure on what information would be relevant to the project and so I
thought the best decision was to record everything. My initial thoughts were to only
take the image portion of the RGBD data; however, after further consideration, the
depth data could also be useful for refining the detection region of the image and
in turn result in a faster model due to a refined search space.
Additionally, my actions could also be used as a form of training data, allowing the
robot to learn from my corrections and inputs. This could be particularly useful in more
tricky situations where basic logic would fail, and would allow for the potential for an
Imitation Learning approach.

Initially, data was extracted as .sov2 files (containing both the image and depth perception
outputs) as well as .csv files (containing the action data). These were timestamped to 
allow for consistency between actions and the vision aspects.

\subsection{Processing}
From here, the .sov2 files were processed to extract independent video and depth data. Video
was chosen as it gave the greatest versatility (and easiest viewing for humans) compared
to individual image extraction. While images are easier to process, the possibility to later
take frames while reducing the visual number of files seemed like the better option---at least
for now.

The depth images were extracted as a compilation of $32 \times 32$ images. These could
then be used to identify the floor (represented as a darker blue section) which would aid in
refining the search space for the yellow line (the path). However, during collection, I
noticed that the shininess of the floor caused a lot of noise in the depth images. This could
make it more difficult to identify the floor and thus may be useless. However, also during
image capture, I noticed that the rope that we are to follow was correctly identified as
being closer (in dark blue), and so it may be possible to use this to help reinforce the
line's location.

\section{Experiments}

\section{Conclusion}

\end{document}